% Compile with pdfLaTeX.
\documentclass[11pt,a4paper]{article}
\usepackage[T1]{fontenc}
\usepackage{lmodern}
\usepackage[margin=25mm]{geometry}
\usepackage{amsmath,amssymb,amsthm}
\usepackage{enumitem}
\usepackage{needspace}
\usepackage[hidelinks]{hyperref}
\usepackage{bookmark}
\hypersetup{pdftitle={Divisorial Proj},pdfauthor={}}
\pagestyle{plain}
\setlength{\parskip}{0.2em}
\newtheorem{lemma}{Lemma}
\newtheorem{proposition}[lemma]{Proposition}
\newtheorem{theorem}[lemma]{Theorem}
\newtheorem{corollary}[lemma]{Corollary}
\theoremstyle{definition}
\newtheorem{remark}[lemma]{Remark}
\theoremstyle{definition}
\newtheorem{example}[lemma]{Example}
\DeclareMathOperator{\Spec}{Spec}
\DeclareMathOperator{\Proj}{Proj}
\DeclareMathOperator{\Supp}{Supp}
\newcommand{\OO}{\mathcal O}
\newcommand{\RR}{\mathcal R}
\setlist[enumerate]{label=\textup{(\arabic*)},leftmargin=2em,itemsep=0.2em}

\begin{document}
\begin{center}
 {\Large\bfseries Divisorial Proj}
\end{center}
\vspace{0.5em}

Let $X$ be a normal integral variety over a field $k$, and let $B$ be an integral Weil divisor on $X$. Assume that the divisorial algebra
\[
 \RR(X,B):=\bigoplus_{m\ge0}\OO_X(mB)
\]
is finitely generated as an $\OO_X$-algebra, and set
\[
 f\colon Z=\Proj_X\RR(X,B)\longrightarrow X.
\]
This note studies the properties of $f$ and proves the following:
\begin{enumerate}
 \item The morphism $f$ detects the non-$\mathbb Q$-Cartier locus of $B$: its isomorphism locus on $X$ is precisely the $\mathbb Q$-Cartier locus of $B$. In particular, $f^{-1}(U)\to U$ is an isomorphism, where $U\subseteq X$ is the smooth locus.
 \item The variety $Z$ is normal.
 \item The morphism $f$ is small. Consequently, the strict transform $B_Z$ of $B$ is $\mathbb Q$-Cartier and $f$-ample.
\end{enumerate}

\setcounter{section}{-1}
\section{The twisting sheaf on Proj}

Let $S=\bigoplus_{n\ge0}S_n$ be a nonnegatively graded ring, and set $P=\Proj S$. For a graded $S$-module $M$, let $\widetilde M$ denote the associated quasicoherent sheaf on $P$, characterized on standard affine opens by
\[
 \Gamma(D_+(s),\widetilde M)=(M_s)_0,
\]
where $s$ is homogeneous of positive degree. Define the shifted graded module $S(1)$ by
\[
 S(1)_n=S_{n+1}.
\]
The twisting sheaf on $P$ is defined as
\[
 \OO_P(1):=\widetilde{S(1)}.
\]
In particular,
\[
 \Gamma(D_+(s),\OO_P(1))
 =(S(1)_s)_0=(S_s)_1.
\]

\begin{lemma}\label{lem:degree-one}
If $S$ is generated by $S_1$ as an $S_0$-algebra, then $\OO_P(1)$ is an invertible sheaf.
\end{lemma}

\begin{proof}
We first show that the standard opens $D_+(s)$, with $s\in S_1$, cover $P$. Indeed, if a homogeneous prime ideal $\mathfrak p$ contains $S_1$, then generation in degree one implies that it contains $S_+$. Such a prime does not belong to $\Proj S$. Hence every point of $P$ lies in $D_+(s)$ for some $s\in S_1$.

Fix such an $s$. Since $s$ is invertible in $S_s$ and has degree one, multiplication by $s$ gives an isomorphism of $(S_s)_0$-modules
\[
 (S_s)_0\xrightarrow{\ \sim\ }(S_s)_1,
 \qquad a\longmapsto sa,
\]
whose inverse is multiplication by $s^{-1}$. Sheafifying on $D_+(s)=\Spec(S_s)_0$ therefore yields
\[
 \OO_{D_+(s)}
 \simeq\OO_P(1)|_{D_+(s)}.
\]
Since these opens cover $P$, the sheaf $\OO_P(1)$ is locally free of rank one.
\end{proof}

The sheaf $\OO_P(1)$ need not be invertible for an arbitrary grading, as the following example shows.

\begin{example}\label{ex:degree-two}
Let $S=k[u]$ with $\deg u=2$. Localizing at $u$ preserves the original grading: $\deg(u^j)=2j$ for every $j\in\mathbb Z$. Thus
\[
 S_u=k[u,u^{-1}]
 =\cdots\oplus ku^{-1}\oplus k\oplus ku\oplus ku^2\oplus\cdots,
\]
where the displayed summands have degrees $\ldots,-2,0,2,4,\ldots$, respectively. In particular, only even degrees occur, and
\[
 (S_u)_0=k,\qquad (S_u)_1=0,\qquad (S_u)_2=ku\simeq k.
\]
Consequently,
\[
 P=D_+(u)=\Spec(S_u)_0\simeq\Spec k,
 \qquad
 \Gamma(P,\OO_P(1))=(S_u)_1=0.
\]
Thus $\OO_P(1)$ is the zero sheaf on the nonempty scheme $P$, and is not invertible.
\end{example}

\section{The isomorphism locus}

Let $X$ be a normal integral variety and $B$ an integral Weil divisor on $X$. Assume that $\RR=\RR(X,B)=\bigoplus_{m\ge0}\OO_X(mB)$ is finitely generated, and let
\[
 f\colon Z=\Proj_X\RR\longrightarrow X.
\]
Choose $r>0$ such that the Veronese algebra
\[
 \RR^{(r)}=\bigoplus_{n\ge0}\OO_X(nrB)
\]
is generated in degree one. Identifying $Z$ with $\Proj_X\RR^{(r)}$, Lemma~\ref{lem:degree-one} shows that
\[
 L:=\OO_{\Proj_X\RR^{(r)}}(1)
\]
is a line bundle on $Z$.

\begin{theorem}\label{prop:isomorphism-locus}
For every point $x\in X$, the following conditions are equivalent:
\begin{enumerate}
 \item $B$ is $\mathbb Q$-Cartier at $x$;
 \item there is an open neighborhood $V$ of $x$ such that $f^{-1}(V)\to V$ is an isomorphism;
 \item the fiber $f^{-1}(x)$ has dimension zero.
\end{enumerate}
\end{theorem}

\begin{proof}
\textup{(1)}$\Rightarrow$\textup{(2)}. Choose a neighborhood $V$ of $x$ and $q>0$ such that $qB|_V$ is Cartier. Then
\[
 f^{-1}(V)
 \simeq\Proj_V\bigoplus_{m\ge0}\OO_V(mqB)
 \simeq\Proj_V\operatorname{Sym}_{\OO_V}\OO_V(qB)
 \simeq V.
\]

\textup{(2)}$\Rightarrow$\textup{(1)}. Identify $f^{-1}(V)$ with $V$. The line bundle $L|_{f^{-1}(V)}$ agrees with $\OO_V(rB)$ on the regular locus of $V$. Both sheaves are reflexive and $V$ is normal, so this isomorphism extends to $V$. Thus $\OO_V(rB)$ is invertible and $B$ is $\mathbb Q$-Cartier at $x$.

The implication \textup{(2)}$\Rightarrow$\textup{(3)} is immediate. For \textup{(3)}$\Rightarrow$\textup{(2)}, since $f$ is projective and birational, upper semicontinuity of fiber dimension gives a neighborhood $V$ of $x$ over which $f$ is quasi-finite~\href{https://stacks.math.columbia.edu/tag/0D4I}{[SP, Tag~0D4I]}. There $f$ is proper and quasi-finite, hence finite~\href{https://stacks.math.columbia.edu/tag/02LS}{[SP, Tag~02LS]}, and therefore an isomorphism because $V$ is normal~\href{https://stacks.math.columbia.edu/tag/0AB1}{[SP, Tag~0AB1]}.
\end{proof}

In particular, if
\[
 \operatorname{Iso}_X(f)
 :=\{x\in X\mid f\text{ is an isomorphism over a neighborhood of }x\},
\]
then Theorem~\ref{prop:isomorphism-locus} identifies this set with the $\mathbb Q$-Cartier locus of $B$. Its complement is
\[
 \begin{aligned}
 X\setminus\operatorname{Iso}_X(f)
 &=\{x\in X\mid B\text{ is not }\mathbb Q\text{-Cartier at }x\}\\
 &=\{x\in X\mid\dim f^{-1}(x)\ge1\}.
 \end{aligned}
\]
Since $\operatorname{Iso}_X(f)$ is open by definition, Theorem~\ref{prop:isomorphism-locus} also shows that the non-$\mathbb Q$-Cartier locus of $B$ is closed.

We give a direct proof of this closedness, independent of Theorem~\ref{prop:isomorphism-locus} and of finite generation. Write
\[
 U_B=\{x\in X\mid B\text{ is Cartier at }x\}.
\]

\begin{lemma}\label{lem:cartier-locus}
The Cartier locus $U_B$ is open in $X$. Equivalently, the non-Cartier locus $X\setminus U_B$ is closed.
\end{lemma}

\begin{proof}
Fix $x\in U_B$. A Cartier divisor on the spectrum of a local ring is principal, since every invertible module over a local ring is free of rank one. Hence there is a rational function $g\in K(X)^*$ such that
\[
 B|_{\Spec\OO_{X,x}}
 =\operatorname{div}_{\Spec\OO_{X,x}}(g).
\]
Set $E=B-\operatorname{div}_X(g)$. The prime divisors of $X$ that contain $x$ correspond to the height-one primes of $\OO_{X,x}$, with the same discrete valuations. The displayed equality therefore says that the coefficient of $E$ along every prime divisor containing $x$ is zero. Thus $x\notin\Supp E$.

The support of a Weil divisor is closed, so $V=X\setminus\Supp E$ is an open neighborhood of $x$. On $V$ we have
\[
 B|_V=\operatorname{div}_V(g),
\]
and consequently $V\subseteq U_B$. This proves that $U_B$ is open.
\end{proof}

\Needspace{12\baselineskip}
\begin{corollary}\label{cor:q-cartier-locus}
The $\mathbb Q$-Cartier locus
\[
 U_B^{\mathbb Q}
 =\{x\in X\mid mB\text{ is Cartier at }x
                 \text{ for some integer }m\ge1\}
\]
is open. In particular, the non-$\mathbb Q$-Cartier locus of $B$ is closed.
\end{corollary}

\begin{proof}
We have $U_B^{\mathbb Q}=\bigcup_{m\ge1}U_{mB}$, and each $U_{mB}$ is open by Lemma~\ref{lem:cartier-locus}.
\end{proof}

\begin{remark}[Locally factorial and $\mathbb Q$-factorial loci]
The analogous openness questions for local factoriality and $\mathbb Q$-factoriality concern all Weil divisors simultaneously. Their loci are the intersections of the Cartier loci, respectively the $\mathbb Q$-Cartier loci, of all Weil divisors. Since these intersections can be infinite, their openness does not follow formally from Lemma~\ref{lem:cartier-locus} and Corollary~\ref{cor:q-cartier-locus}.

For a normal variety over an algebraically closed field, both loci are open, by Boissi\`ere--Gabber--Serman~\cite[Theorems~4.2 and~6.1]{BGS}; hence their complements are closed. Here a variety is an integral separated scheme of finite type over the ground field, and these results hold in arbitrary characteristic. The non-$\mathbb Q$-Cartier locus of a fixed divisor $B$ is contained in the non-$\mathbb Q$-factorial locus, but can be strictly smaller: for instance, $B=0$ is Cartier everywhere.
\end{remark}

\Needspace{10\baselineskip}
\section{Normality of the divisorial Proj}

We use the isomorphism over the regular locus to prove that $Z$ is normal.

Let $X$ be a normal integral variety and $B$ an integral Weil divisor such that $\RR(X,B)=\bigoplus_{m\ge0}\OO_X(mB)$ is finitely generated. Set
\[
 f\colon Z=\Proj_X\RR(X,B)\longrightarrow X.
\]
Choose $r>0$ such that $\RR(X,B)^{(r)}$ is generated in degree one, and let
\[
 L:=\OO_{\Proj_X\RR(X,B)^{(r)}}(1)
\]
be the tautological line bundle on $Z\simeq\Proj_X\RR(X,B)^{(r)}$.

\begin{theorem}\label{prop:normality}
The variety $Z$ is normal, and the canonical maps give isomorphisms
\[
 \OO_X(mrB)\xrightarrow{\ \sim\ }f_*L^{\otimes m}
 \qquad(m\ge0).
\]
\end{theorem}

\begin{proof}
Let $\nu\colon W\to Z$ be the normalization, and set
\[
 g=f\circ\nu\colon W\to X,
 \qquad L'=\nu^*L.
\]
Since $Z$ is a variety, $\nu$ is finite~\href{https://stacks.math.columbia.edu/tag/0BXR}{[SP, Tag~0BXR]}. Hence $g$ is proper, and $L'$ is $g$-ample: over each affine open subset of $X$, it is the pullback of an ample line bundle by a finite morphism.

Let $U=X_{\mathrm{reg}}$ and $j\colon U\hookrightarrow X$ be the inclusion. Over an algebraically closed field, $U$ is the smooth locus. Since $X$ is normal, $X\setminus U$ has codimension at least two. Every Weil divisor on $U$ is Cartier, so Theorem~\ref{prop:isomorphism-locus} identifies $f^{-1}(U)$ with $U$. The normalization is also an isomorphism there. Under these identifications,
\[
 L|_{f^{-1}(U)}\simeq\OO_U(rB),
 \qquad
 L'|_{g^{-1}(U)}\simeq\OO_U(rB).
\]

Fix $m\ge0$. Restriction to the dense open subset $g^{-1}(U)$ gives an injection
\[
 g_*(L')^{\otimes m}
 \hookrightarrow j_*\OO_U(mrB)
 =\OO_X(mrB).
\]
The injection holds because $W$ is integral and $(L')^{\otimes m}$ is invertible; the last equality follows from reflexivity of $\OO_X(mrB)$ and the codimension of $X\setminus U$. In the other direction, the defining graded algebra and pullback by $\nu$ give natural maps
\[
 \OO_X(mrB)\longrightarrow f_*L^{\otimes m}
 \longrightarrow g_*(L')^{\otimes m}.
\]
The second map is injective because $\nu$ is birational and $Z$ is integral. The composite of these maps with the preceding injection into $\OO_X(mrB)$ is the identity: each section restricts to the same rational function on $U$. It follows that all the maps are isomorphisms. In particular,
\begin{equation}\label{eq:normalization-section-algebra}
 \bigoplus_{m\ge0}g_*(L')^{\otimes m}
 \simeq\bigoplus_{m\ge0}\OO_X(mrB)
 =\RR^{(r)}
\end{equation}
as graded $\OO_X$-algebras; the identifications respect multiplication because they are induced by restriction to $U$.

A proper morphism equipped with a relatively ample line bundle is recovered as the relative Proj of its section algebra~\href{https://stacks.math.columbia.edu/tag/0C6J}{[SP, Tag~0C6J]}. Applying this to $(g,L')$ and using~\eqref{eq:normalization-section-algebra}, we obtain
\[
 W\xrightarrow{\ \sim\ }
 \Proj_X\bigoplus_{m\ge0}g_*(L')^{\otimes m}
 \simeq\Proj_X\RR^{(r)}=Z.
\]
The resulting morphism $W\to Z$ is $\nu$, since it is induced by the pullbacks of the defining homogeneous sections. Thus $\nu$ is an isomorphism and $Z$ is normal. The asserted direct-image identities were established by the same comparison.
\end{proof}

\Needspace{12\baselineskip}
\section{Smallness}

To prepare for the proof of smallness, we express $f$ locally as a blowup.

On an affine open subset $V\subseteq X$, choose an effective Cartier divisor $H$ such that
\[
 J:=\OO_V(rB-H)\subseteq\OO_V
\]
is an ideal sheaf. Since $\RR^{(r)}$ is generated in degree one,
\[
 \OO_V(rB)^m=\OO_V(mrB)\qquad(m\ge0).
\]
Hence
\[
 J^m=\OO_V(rB)^m\otimes\OO_V(-mH)
     =\OO_V(mrB-mH),
\]
and
\[
 \bigoplus_{m\ge0}J^m
 =\bigoplus_{m\ge0}\OO_V(mrB)\otimes\OO_V(-H)^{\otimes m}.
\]
Twisting by a line bundle does not change relative Proj~\href{https://stacks.math.columbia.edu/tag/02NC}{[SP, Tag~02NC]}, so
\begin{equation}\label{eq:blowup-presentation}
 \operatorname{Bl}_J V
 =\Proj_V\bigoplus_{m\ge0}J^m
 \simeq\Proj_V(\RR^{(r)}|_V)
 \simeq f^{-1}(V).
\end{equation}
By Theorem~\ref{prop:isomorphism-locus}, its non-isomorphism locus is the non-$\mathbb Q$-Cartier locus of $B|_V$.

\Needspace{10\baselineskip}
\begin{proposition}\label{prop:smallness}
The morphism $f$ is small. If $B_Z$ denotes the strict transform of $B$, then
\[
 L\simeq\OO_Z(rB_Z).
\]
In particular, $B_Z$ is $\mathbb Q$-Cartier and $f$-ample.
\end{proposition}

\begin{proof}
Smallness is local on $X$, so work over an affine open subset and use the blowup presentation above. Set
\[
 D=H-rB\ge0,\qquad J=\OO_X(-D),
\]
and define the effective Cartier divisor $F$ by
\[
 J\OO_Z=\OO_Z(-F).
\]
Then $J^m=\OO_X(-mD)$ for all $m\ge0$, and $\OO_Z(-F)$ is $f$-ample. On the big open set $U=X_{\mathrm{reg}}$, the morphism $f$ is an isomorphism and $F$ agrees with $D$. Thus $(f_*F)|_U=D|_U$, and hence
\[
 f_*F=D.
\]

Suppose that $f$ contracts a prime divisor $E$. Since $Z$ is normal and $\OO_Z(-F)$ is ample over the affine base, $\OO_Z(-mF+E)$ is globally generated for $m\gg0$~\href{https://stacks.math.columbia.edu/tag/01Q3}{[SP, Tag~01Q3]}. Choose a section that does not vanish at $\eta_E$ and represent it by $\varphi\in K(X)^*$. Then
\[
 C:=\operatorname{div}_Z(\varphi)-mF+E\ge0,
 \qquad \operatorname{coeff}_E C=0.
\]
Consequently,
\begin{equation}\label{eq:smallness-negative-coefficient}
 \operatorname{coeff}_E\bigl(\operatorname{div}_Z(\varphi)-mF\bigr)=-1.
\end{equation}
Pushing $C$ down to $X$ and using $f_*E=0$ gives
\[
 \operatorname{div}_X(\varphi)-mD=f_*C\ge0.
\]
Therefore
\[
 \varphi\in H^0(X,\OO_X(-mD))=H^0(X,J^m).
\]
Since $J^m\OO_Z=\OO_Z(-mF)$, it follows that
\[
 \operatorname{div}_Z(\varphi)-mF\ge0,
\]
contradicting~\eqref{eq:smallness-negative-coefficient}. Thus $f$ is small.

It follows that $F=D_Z$ and $f^*H=H_Z$, where subscripts denote strict transforms. The formula for the tautological line bundle under twisting now gives
\[
 L\simeq\OO_Z(f^*H-F)
 =\OO_Z(H_Z-D_Z)
 =\OO_Z(rB_Z).
\]
These local identifications are induced by the rational realizations of the sheaves and hence glue. Since $L$ is $f$-ample, $B_Z$ is $\mathbb Q$-Cartier and $f$-ample.
\end{proof}

\begin{thebibliography}{BGS}
\bibitem[SP]{Stacks}
The Stacks Project Authors,
\emph{The Stacks Project},
\href{https://stacks.math.columbia.edu/tag/0EGH}{Lemma~10.56.2}
and Tags
\href{https://stacks.math.columbia.edu/tag/0B5J}{0B5J},
\href{https://stacks.math.columbia.edu/tag/02NC}{02NC},
\href{https://stacks.math.columbia.edu/tag/01OF}{01OF},
\href{https://stacks.math.columbia.edu/tag/01Q3}{01Q3},
\href{https://stacks.math.columbia.edu/tag/0D4I}{0D4I},
\href{https://stacks.math.columbia.edu/tag/02LS}{02LS},
\href{https://stacks.math.columbia.edu/tag/0AB1}{0AB1},
\href{https://stacks.math.columbia.edu/tag/0BXR}{0BXR},
\href{https://stacks.math.columbia.edu/tag/0C6J}{0C6J}.

\bibitem{BGS}
S.~Boissi\`ere, O.~Gabber, and O.~Serman,
\emph{Sur le produit de vari\'et\'es localement factorielles ou $\mathbb Q$-factorielles},
\href{https://arxiv.org/abs/1104.1861}{arXiv:1104.1861}, version~4 (2019).
\end{thebibliography}

\end{document}
